% Additional lessons: trees, graphs, and dynamic programming.

\patternintoc{Tree Traversal and Recursion}
\problemstart{Binary Tree Level Order Traversal}{102}{Medium}{Breadth-first search}{Queue}

\subsection{The Problem}
Given the root of a binary tree, return its values one level at a time from
top to bottom. Nodes on the same depth belong to the same output row.

\begin{conceptbox}{Example}
For the tree with root \(3\), children \(9\) and \(20\), and children
\(15\) and \(7\) under \(20\), return
\texttt{[[3],[9,20],[15,7]]}. An empty tree returns an empty vector.
\end{conceptbox}

\subsection{From the Direct Idea to BFS}
A depth-first traversal can record each node's depth, but it needs extra logic
to place values into the correct row. The cleaner approach processes the tree in
the same order as the required output: breadth first. A queue stores the next
nodes to visit.

At the beginning of each outer-loop iteration, the queue contains exactly one
complete level. Save the current queue size before adding any children. That
size tells us how many nodes belong to the current row.

\begin{visualbox}{The queue handles one tree level at a time}
\centering
\begin{tikzpicture}[>={Stealth[length=2.2mm,width=1.55mm]},font=\small,level distance=12mm]
  \node[trienode,visited] (r) {3};
  \node[trienode,below left=12mm and 18mm of r] (a) {9};
  \node[trienode,below right=12mm and 18mm of r] (b) {20};
  \node[trienode,below left=12mm and 8mm of b] (c) {15};
  \node[trienode,below right=12mm and 8mm of b] (d) {7};
  \draw[flowarrow] (r)--(a); \draw[flowarrow] (r)--(b);
  \draw[flowarrow] (b)--(c); \draw[flowarrow] (b)--(d);
  \node[draw=teal,dashed,fit=(r),inner sep=3mm,label=right:{\scriptsize level 0}] {};
  \node[draw=orange,dashed,fit=(a)(b),inner sep=3mm,label=right:{\scriptsize level 1}] {};
  \node[draw=teal,dashed,fit=(c)(d),inner sep=3mm,label=right:{\scriptsize level 2}] {};
\end{tikzpicture}

\smallskip
Queue after each level: \texttt{[3]} $\rightarrow$ \texttt{[9,20]}
$\rightarrow$ \texttt{[15,7]}.
\end{visualbox}

\subsection{Algorithm and Key Idea}
Start by pushing the root. While the queue is not empty, save
\texttt{levelSize}, remove exactly that many nodes, append their values to one
row, and push each non-null child. The key rule is that every queued node is
unprocessed and the first \texttt{levelSize} nodes share the same depth.

\begin{invariantbox}{one saved queue prefix equals one tree level}
Immediately before an outer-loop round, the queue contains all unprocessed
nodes in breadth-first order. The saved size identifies exactly the current
depth; children appended during the round belong to the next depth and are
so not consumed until the following round.
\end{invariantbox}

\subsection{Pseudocode}
\begin{lstlisting}[style=pseudocode]
levelOrder(root):
    if root is null: return []
    queue = [root]
    result = []
    while queue is not empty:
        levelSize = queue.size()
        level = []
        repeat levelSize times:
            node = pop front
            append node.value to level
            push non-null children
        append level to result
    return result
\end{lstlisting}

\subsection{C++ Implementation}
\begin{lstlisting}[style=compactcode,caption={Complete C++ solution: Binary Tree Level Order Traversal}]
#include <queue>
#include <utility>
#include <vector>

using namespace std;

class Solution {
public:
    vector<vector<int>> levelOrder(TreeNode* root) {
        if (root == nullptr) return {};

        vector<vector<int>> result;
        queue<TreeNode*> pending;
        pending.push(root);

        while (!pending.empty()) {
            const int levelSize = static_cast<int>(pending.size());
            vector<int> level;
            level.reserve(levelSize);

            for (int i = 0; i < levelSize; ++i) {
                TreeNode* node = pending.front();
                pending.pop();
                level.push_back(node->val);
                if (node->left != nullptr) pending.push(node->left);
                if (node->right != nullptr) pending.push(node->right);
            }
            result.push_back(move(level));
        }
        return result;
    }
};
\end{lstlisting}

\subsection{Time, Space, and Common Mistakes}
Every node enters and leaves the queue once, so time is \(O(n)\). The queue
uses \(O(w)\) space, where \(w\) is the maximum tree width; the returned rows
use \(O(n)\) output space. Save the level size before pushing children, and
handle a null root before accessing it.

\begin{revisionbox}{How to explain it in an interview}
I use BFS because the output is grouped by depth. The queue contains the next
level. I snapshot its size, process exactly those nodes, and enqueue their
children for the following level. Each node is processed once: \(O(n)\) time
and \(O(w)\) extra queue space.
\end{revisionbox}

\subsection{Detailed Visual Dry Run}
For the worked example, the queue changes by complete levels. Children are
added only after the current node is removed.

\begin{dryrunbox}{Queue trace for \texttt{[3,9,20,null,null,15,7]}}
\begin{center}
\begin{tabularx}{\linewidth}{@{}c c X c@{}}
\toprule
Round & Saved size & Nodes processed and children appended & Output \\
\midrule
1 & 1 & Remove 3; append 9 and 20 & \texttt{[[3]]} \\
2 & 2 & Remove 9; remove 20 and append 15 and 7 & \texttt{[[3],[9,20]]} \\
3 & 2 & Remove 15 and 7; neither has children & \texttt{[[3],[9,20],[15,7]]} \\
\bottomrule
\end{tabularx}
\end{center}
\end{dryrunbox}

\begin{visualbox}{What the queue holds after each level}
\centering
\begin{tikzpicture}[font=\small,>={Stealth[length=2.2mm,width=1.55mm]}]
  \node[text=teal,font=\bfseries] at (0.2,2.05) {front on the left};
  \node[text=orange,font=\bfseries] at (4.3,2.05) {back on the right};

  \node[text=navy,font=\bfseries,anchor=east] at (-0.8,1.3) {start};
  \node[arraycell,minimum width=13mm] (q0) at (0,1.3) {3};
  \node[right=12mm of q0,text=navy,font=\scriptsize] {nodes in this level \(=1\)};

  \node[text=navy,font=\bfseries,anchor=east] at (-0.8,0) {after level 0};
  \node[arraycell,minimum width=13mm] (q1) at (0,0) {9};
  \node[arraycell,minimum width=13mm,right=0mm of q1] (q2) {20};
  \node[right=12mm of q2,text=navy,font=\scriptsize] {nodes in this level \(=2\)};

  \node[text=navy,font=\bfseries,anchor=east] at (-0.8,-1.3) {after level 1};
  \node[arraycell,minimum width=13mm] (q3) at (0,-1.3) {15};
  \node[arraycell,minimum width=13mm,right=0mm of q3] (q4) {7};
  \node[right=12mm of q4,text=navy,font=\scriptsize] {nodes in this level \(=2\)};
\end{tikzpicture}
\end{visualbox}

The saved size is essential. If the loop instead continued until the queue
became empty, the children would be consumed in the same row as their parents
and the level grouping would be lost.

\subsection{Code Walkthrough}
\begin{enumerate}
  \item Reject a null root and return an empty result.
  \item Push the root into \texttt{pending}. The queue now represents the
  first unprocessed level.
  \item Copy \texttt{pending.size()} into \texttt{levelSize}. This value does
  not change when children are pushed later in the round.
  \item Pop exactly \texttt{levelSize} nodes, record their values, and enqueue
  their non-null children from left to right.
  \item Move the completed row into the result and repeat for the next depth.
\end{enumerate}

Although the code contains a loop inside another loop, it is still linear.
The inner loops across all rounds process each node exactly once; they do not
restart a full scan for every level.

\subsection{Edge Cases, Follow-ups, and Practice}
Test an empty tree, one node, a completely skewed tree, and a wide final
level. A common mistake is to enqueue null children and then dereference them.
Useful follow-ups include zigzag level order, returning only the rightmost
node at each depth, and computing the average value of every level.

\begin{revisionbox}{Practice checkpoint}
Without looking at the code, explain what the queue contains immediately
before each outer-loop iteration. Then modify the algorithm so alternating
rows are reversed while the traversal remains breadth first.
\end{revisionbox}
\chaptercheckpoint
  {The requested output is grouped by tree depth.}
  {A FIFO queue and the size of the current level.}
  {Processing newly enqueued children in the current output row.}
  {$O(n)$ time and $O(w)$ extra queue space.}
\judgeclassnote


\problemstart{Binary Tree Maximum Path Sum}{124}{Hard}{Postorder DFS}{Binary tree}

\subsection{The Problem}
A path may start and end at any nodes, but it cannot reuse a node. Return the
largest sum of any non-empty path. The path does not have to pass through the
root.

\begin{conceptbox}{Example}
For \texttt{[-10,9,20,null,null,15,7]}, the best path is
\(15\rightarrow20\rightarrow7\), so the answer is \(42\).
\end{conceptbox}

\subsection{Building the Recursive State}
A brute-force search could start a path at every node and explore many of the
same subtrees repeatedly. Postorder DFS removes that repetition. A parent can
continue through at most one child, so each recursive call returns the best
one-sided gain beginning at its node.

For a node with value \(x\): \(L=\max(0,\text{left gain}),\qquad R=\max(0,\text{right gain}).\)
The complete path turning at this node is \(L+x+R\). The value returned to the
parent is \(x+\max(L,R)\).

\begin{invariantbox}{a returned gain is always a single extendable chain}
After \texttt{gain(node)} finishes, its return value is the maximum sum of a
path that starts at \texttt{node} and travels downward through at most one
child. The global answer separately stores the best complete path seen so far,
which may use both children at exactly one turning node.
\end{invariantbox}

\begin{visualbox}{A path can join both children at one middle node}
\centering
\begin{tikzpicture}[>={Stealth[length=2.2mm,width=1.55mm]},font=\small]
  \node[trienode,active] (r) {20};
  \node[trienode,visited,below left=14mm and 16mm of r] (l) {15};
  \node[trienode,visited,below right=14mm and 16mm of r] (q) {7};
  \draw[-,very thick,orange] (l)--(r)--(q);
  \node[below=9mm of l,text=teal] {best left path: 15};
  \node[below=9mm of q,text=teal] {best right path: 7};
  \node[above=7mm of r,draw=orange,fill=softorange,rounded corners]
    {full path: \(15+20+7=42\)};
\end{tikzpicture}
\end{visualbox}

\begin{visualbox}{A node sends only one branch up to its parent}
\centering
\begin{tikzpicture}[font=\small,>={Stealth[length=2.2mm,width=1.55mm]}]
  \node[trienode,active] (r) {20};
  \node[trienode,visited,below left=13mm and 16mm of r] (l) {15};
  \node[trienode,visited,below right=13mm and 16mm of r] (q) {7};
  \draw[-,thick,navy] (r)--(l) (r)--(q);
  \node[left=7mm of l,align=right,text=teal] {sends up\\15};
  \node[right=7mm of q,align=left,text=teal] {sends up\\7};
  \node[dsnode,right=35mm of r,fill=softorange] (full)
    {best full path\\\(15+20+7=42\)};
  \node[dsnode,below=9mm of full,fill=softteal] (up)
    {send up to parent\\\(20+\max(15,7)=35\)};
  \draw[flowarrow,orange] (r.east) -- (full.west);
  \draw[flowarrow,teal] (full) -- (up);
\end{tikzpicture}
\end{visualbox}

\subsection{Pseudocode}
\begin{lstlisting}[style=pseudocode]
dfs(node):
    if node is null: return 0
    leftGain = max(0, dfs(node.left))
    rightGain = max(0, dfs(node.right))
    answer = max(answer, leftGain + node.value + rightGain)
    return node.value + max(leftGain, rightGain)
\end{lstlisting}

\subsection{C++ Implementation}
\begin{lstlisting}[style=compactcode,caption={Complete C++ solution: Binary Tree Maximum Path Sum}]
#include <algorithm>
#include <limits>

using namespace std;

class Solution {
    int bestPath = numeric_limits<int>::min();

    int gain(TreeNode* node) {
        if (node == nullptr) return 0;
        const int left = max(0, gain(node->left));
        const int right = max(0, gain(node->right));
        bestPath = max(bestPath, node->val + left + right);
        return node->val + max(left, right);
    }

public:
    int maxPathSum(TreeNode* root) {
        bestPath = numeric_limits<int>::min();
        gain(root);
        return bestPath;
    }
};
\end{lstlisting}

\subsection{Why It Works and Complexity}
Postorder traversal makes both child gains available before processing the
parent. Every possible path has a highest turning node; the algorithm tests
that path when it processes that node. Negative gains are discarded because
they can only reduce the sum. Time is \(O(n)\), and recursion uses \(O(h)\)
stack space for tree height \(h\).

\begin{revisionbox}{Common mistakes}
Do not start the answer to zero because all node values may be negative.
Do not return the two-child path to the parent; that would branch and stop
being a valid path. Return only the better one-sided gain.
\end{revisionbox}

\begin{conceptbox}{All-negative tree}
For a tree such as \texttt{[-3,-2,-5]}, every child gain is clipped to zero,
but each node still tests its own value as a non-empty path. Initializing the
global answer to negative infinity returns \(-2\), while
initializing it to zero would incorrectly allow an empty path.
\end{conceptbox}

\subsection{From Brute Force to One Postorder Traversal}
A direct method could treat every node as a possible path start and explore
every possible continuation. Long chains and overlapping subtrees would then
be traversed repeatedly, leading to \(O(n^2)\) work on a skewed tree. The
optimized method asks each subtree one reusable question: ``What is the best
one-sided gain that can be extended through your parent?''

This returned value and the global answer are deliberately different:
\begin{itemize}
  \item the value returned upward may use at most one child;
  \item the complete choice recorded locally may join both child gains.
\end{itemize}

\subsection{Detailed Visual Dry Run}
\begin{dryrunbox}{Postorder values for \texttt{[-10,9,20,null,null,15,7]}}
\begin{center}
\begin{tabular}{@{}c c c c c@{}}
\toprule
Node & Left gain & Right gain & Complete choice & Gain returned \\
\midrule
9 & 0 & 0 & 9 & 9 \\
15 & 0 & 0 & 15 & 15 \\
7 & 0 & 0 & 7 & 7 \\
20 & 15 & 7 & 42 & 35 \\
-10 & 9 & 35 & 34 & 25 \\
\bottomrule
\end{tabular}
\end{center}
The best complete path is recorded at node 20. The value 42 is not returned to
node \(-10\), because returning both branches would create an invalid fork.
\end{dryrunbox}

\subsection{Code Walkthrough}
The helper first recurses left and right, so it is a postorder traversal.
\texttt{max(0, gain)} discards a negative branch. The expression
\texttt{node->val + left + right} tests the path whose highest turning point
is the current node. The helper then returns the current value plus only the
larger child gain. Resetting \texttt{bestPath} inside the public method makes
the same \texttt{Solution} object safe for another call.

\subsection{Edge Cases and Interview Follow-ups}
The tree may contain one negative node, all negative nodes, or a best path
that never touches the root. Recursion uses \(O(h)\) call-stack space and can
reach \(O(n)\) for a skewed tree. A useful follow-up is to return the actual
nodes in the maximum path; that requires storing or reconstructing the chosen
child direction as well as the sum.

\begin{revisionbox}{How to explain it simply}
Each recursive call returns the best downward path the parent can extend. At
each node I also test the complete path that joins the best non-negative left
and right gains. Every path has exactly one highest turning node, so one of
these local choices represents the global optimum.
\end{revisionbox}
\chaptercheckpoint
  {A path may begin and end anywhere, and a parent needs child results.}
  {One-sided postorder gains plus a global complete-path maximum.}
  {Returning the two-child choice upward or allowing an empty path.}
  {$O(n)$ time and $O(h)$ recursion space.}
\judgeclassnote


\patternintoc{Graph Traversal}
\problemstart{Clone Graph}{133}{Medium}{DFS with memoization}{Graph and hash map}

\subsection{The Problem}
Given one node in a connected undirected graph, return a deep copy. Every
original node must correspond to one new node, and neighbor relationships must
be keepd without sharing original pointers.

\subsection{Why a Hash Map Is Needed}
A plain recursive copy loops forever on cycles and may clone the same node more
than once. Store \(\text{original pointer}\rightarrow\text{clone pointer}\)
before recursing into neighbors. Seeing an existing key means the clone
is already available.

\begin{invariantbox}{every discovered original has exactly one registered clone}
Before DFS follows an outgoing edge, the current original node already has a
map entry. As a result, a back edge, self-loop, or shared neighbor always
returns the same clone pointer instead of allocating another object.
\end{invariantbox}

\begin{visualbox}{Copy every node and every connection}
\centering
\begin{tikzpicture}[>={Stealth[length=2.2mm,width=1.55mm]},font=\small]
  \node[trienode] (a) {1}; \node[trienode,right=16mm of a] (b) {2};
  \node[trienode,below=14mm of b] (c) {3}; \node[trienode,below=14mm of a] (d) {4};
  \draw[-,navy,thick] (a)--(b)--(c)--(d)--(a);
  \node[above=7mm of a,text=navy,font=\bfseries] {original};
  \node[right=18mm of b,font=\Large,text=teal] (arrow) {$\Longrightarrow$};
  \node[trienode,visited,right=18mm of arrow] (aa) {1'};
  \node[trienode,visited,right=16mm of aa] (bb) {2'};
  \node[trienode,visited,below=14mm of bb] (cc) {3'};
  \node[trienode,visited,below=14mm of aa] (dd) {4'};
  \draw[-,teal,thick] (aa)--(bb)--(cc)--(dd)--(aa);
  \node[above=7mm of aa,text=teal,font=\bfseries] {copy};
\end{tikzpicture}
\end{visualbox}

\subsection{Pseudocode}
\begin{lstlisting}[style=pseudocode]
clone(node):
    if node is null: return null
    if node is in cloneMap: return cloneMap[node]
    copy = new Node(node.value)
    cloneMap[node] = copy
    for neighbor in node.neighbors:
        copy.neighbors.push_back(clone(neighbor))
    return copy
\end{lstlisting}

\subsection{C++ Implementation}
\begin{lstlisting}[style=compactcode,caption={Complete C++ solution: Clone Graph}]
#include <unordered_map>

using namespace std;

class Solution {
    unordered_map<Node*, Node*> clones;

    Node* clone(Node* node) {
        if (node == nullptr) return nullptr;
        const auto found = clones.find(node);
        if (found != clones.end()) return found->second;

        Node* copy = new Node(node->val);
        clones[node] = copy;
        for (Node* neighbor : node->neighbors) {
            copy->neighbors.push_back(clone(neighbor));
        }
        return copy;
    }

public:
    Node* cloneGraph(Node* node) {
        clones.clear();
        return clone(node);
    }
};
\end{lstlisting}

\subsection{Why It Works, Time, and Space}
The map creates exactly one clone for every reachable original node. Each edge
is copied when its adjacency entry is visited. Time is \(O(V+E)\), map and
recursion space are \(O(V)\). Insert the clone into the map before visiting
neighbors; doing it afterward fails on cycles and self-loops.

\begin{revisionbox}{How to explain it in an interview}
I run DFS and memoize original-to-copy pointers. If a node was cloned already,
I reuse it. Otherwise I create and register the clone first, then recursively
copy every neighbor. This handles cycles and keeps shared neighbors.
\end{revisionbox}

\subsection{Why a Plain Recursive Copy Fails}
On a tree, recursively copying children eventually reaches null. A graph may
return to a node already on the active recursion path. Without the map, the
copy of node 1 asks for node 2, node 2 asks for node 1 again, and recursion
never ends. The map is both a cycle guard and the record that
keeps one-to-one node identity.

\subsection{Detailed Map and DFS Trace}
\begin{dryrunbox}{Cloning the four-node cycle}
\begin{center}
\begin{tabularx}{\linewidth}{@{}c X X@{}}
\toprule
Visit & Map before the recursive calls & Action \\
\midrule
1 & empty & create \(1'\), store \(1\to1'\), then visit its neighbors \\
2 & \(1\to1'\) & create \(2'\), store it, then visit neighbor 1 \\
1 again & \(1\to1',2\to2'\) & reuse \(1'\); do not recurse again \\
3 and 4 & earlier clones remain stored & create each once and copy every edge \\
\bottomrule
\end{tabularx}
\end{center}
\end{dryrunbox}

\begin{visualbox}{Each original node points to its one copy}
\centering
\begin{tikzpicture}[font=\small,>={Stealth[length=2.2mm,width=1.55mm]}]
  \node[dsnode,fill=softblue] (o1) {old node 1};
  \node[dsnode,below=7mm of o1,fill=softblue] (o2) {old node 2};
  \node[dsnode,below=7mm of o2,fill=softblue] (o3) {old node 3};
  \node[dsnode,right=38mm of o1,fill=softteal] (c1) {copy \(1'\)};
  \node[dsnode,below=7mm of c1,fill=softteal] (c2) {copy \(2'\)};
  \node[dsnode,below=7mm of c2,fill=softteal] (c3) {copy \(3'\)};
  \draw[flowarrow,teal] (o1)--node[edgelabel,above,font=\scriptsize]{saved pair} (c1);
  \draw[flowarrow,teal] (o2)--(c2);
  \draw[flowarrow,teal] (o3)--(c3);
  \node[draw=navy,dashed,fit=(o1)(o2)(o3)(c1)(c2)(c3),inner sep=5mm,
    label=above:{\small\textbf{saved old-to-copy pairs}}] {};
\end{tikzpicture}

\smallskip
If an old node appears again, the saved pair returns the same copy. This keeps
all shared connections and cycles correct.
\end{visualbox}

\begin{visualbox}{If a node was copied already, stop going deeper}
\centering
\begin{tikzpicture}[font=\small,>={Stealth[length=2.2mm,width=1.55mm]}]
  \node[draw=navy,fill=softblue,minimum width=43mm,minimum height=8mm] (g0)
    {copy node 1};
  \node[draw=navy,fill=softblue,minimum width=43mm,minimum height=8mm,
    above=0mm of g0] (g1) {copy node 2};
  \node[draw=orange,fill=softorange,minimum width=43mm,minimum height=8mm,
    above=0mm of g1] (g2) {see node 1 again};
  \node[left=12mm of g2,text=orange,font=\bfseries,align=right] (top) {working\\now};
  \draw[flowarrow,orange] (top)--(g2);
  \node[dsnode,right=25mm of g1,fill=softteal] (hit)
    {node 1 already has copy \(1'\)\\use that copy};
  \draw[flowarrow,teal] (g2.east)--(hit.north west);
\end{tikzpicture}

\smallskip
The repeated node creates nothing new. Use its saved copy and finish this
function call.
\end{visualbox}

\subsection{Code Walkthrough}
\begin{enumerate}
  \item A null input represents an empty graph and returns null.
  \item Search \texttt{clones} before allocating. An existing entry means the
  node and all references to its identity already have a copy.
  \item Allocate a new node and immediately place it in the map.
  \item Recursively clone every neighbor and append the returned pointer to
  the copy's adjacency vector in the same order.
\end{enumerate}

The immediate insertion in step 3 is the key rule-preserving step. It makes
the partial clone visible before an edge can lead back to it.

\subsection{Edge Cases, Alternative, and Practice}
Test a single isolated node, a self-loop, two nodes connected to each other,
and several nodes sharing the same neighbor. Verify that cloned pointers are
different from all original pointers. An iterative BFS alternative uses a
queue plus the same original-to-clone map and has the same \(O(V+E)\) cost.

\begin{revisionbox}{Practice checkpoint}
Draw a graph containing a self-loop and explain exactly why registering the
copy after cloning its neighbors would recurse forever.
\end{revisionbox}
\chaptercheckpoint
  {A pointer-based graph must be copied without sharing original objects.}
  {An original-pointer-to-clone-pointer hash map and DFS stack.}
  {Registering a clone only after recursively visiting its neighbors.}
  {$O(V+E)$ time and $O(V)$ map plus traversal space.}
\judgeclassnote


\patternintoc{Dynamic Programming}
\problemstart{Coin Change}{322}{Medium}{Bottom-up dynamic programming}{One-dimensional DP}

\subsection{The Problem}
Given coin denominations and a target amount, return the fewest coins needed
to form that amount. Coins may be reused. Return \(-1\) when the amount is
unreachable.

\begin{conceptbox}{Example}
With \texttt{coins=[1,2,5]} and \texttt{amount=11}, the best construction is
\(5+5+1\), so the answer is \(3\).
\end{conceptbox}

\subsection{How We Get to the Solution}
Trying every sequence of coins creates an exponential recursion tree. The same
remaining amounts occur repeatedly. Let \(dp[x]\) be the minimum number of
coins needed to form \(x\). Start \(dp[0]=0\) and every other entry to an
unreachable sentinel \(\texttt{amount}+1\).

For each amount \(x\) and each coin \(c\le x\): \(dp[x]=\min(dp[x],1+dp[x-c]).\)

\begin{invariantbox}{every smaller amount is already best}
When the outer loop reaches amount \(x\), each referenced state \(dp[x-c]\)
has already been finalized. Trying every possible final coin tests
every construction of \(x\), and the minimum choice becomes best.
\end{invariantbox}

\begin{visualbox}{To make 11, add coin 5 to the best answer for 6}
\centering
\begin{tikzpicture}[font=\small]
  \foreach \x/\v in {0/0,1/1,2/1,3/2,4/2,5/1,6/2,7/2,8/3,9/3,10/2,11/3} {
    \node[arraycell] (d\x) at (0.8*\x,0) {\v};
    \node[below=1mm of d\x,font=\scriptsize] {\x};
  }
  \draw[flowarrow,orange] (d6.north) to[bend left=28]
    node[edgelabel,above,font=\scriptsize]{add coin 5} (d11.north);
\end{tikzpicture}
\end{visualbox}

\subsection{Pseudocode}
\begin{lstlisting}[style=pseudocode]
dp = array(amount + 1, amount + 1)
dp[0] = 0
for current from 1 through amount:
    for coin in coins:
        if coin <= current:
            dp[current] = min(dp[current], 1 + dp[current - coin])
return -1 if dp[amount] is unreachable, otherwise dp[amount]
\end{lstlisting}

\subsection{C++ Implementation}
\begin{lstlisting}[style=compactcode,caption={Complete C++ solution: Coin Change}]
#include <algorithm>
#include <vector>

using namespace std;

class Solution {
public:
    int coinChange(vector<int>& coins, int amount) {
        vector<int> dp(amount + 1, amount + 1);
        dp[0] = 0;

        for (int current = 1; current <= amount; ++current) {
            for (const int coin : coins) {
                if (coin <= current) {
                    dp[current] = min(dp[current], 1 + dp[current - coin]);
                }
            }
        }
        return dp[amount] == amount + 1 ? -1 : dp[amount];
    }
};
\end{lstlisting}

\subsection{Time and Space}
For \(A=\texttt{amount}\) and \(m\) coin types, time is \(O(A m)\) and space
is \(O(A)\). The key rule is that before computing \(dp[x]\), every smaller
amount is already best. Avoid using zero as the default for unreachable
amounts; it would look like a free solution.

\begin{revisionbox}{Pattern clue}
When a target can be built from reusable choices and the question asks for a
minimum, maximum, or count, define the answer for every smaller target and
extend those solved states.
\end{revisionbox}

\subsection{Brute Force and Repeated Work}
A recursive brute-force solution tries every coin, subtracts it from the
remaining amount, and repeats. For amount 11, branches such as
\(11\to10\to8\) and \(11\to9\to8\) both solve the same remaining amount 8.
Without memoization, the recursion tree grows exponentially. The DP table
stores the answer to each remaining amount once.

\subsection{Detailed Table Walkthrough}
\begin{dryrunbox}{Building \(dp\) for coins \texttt{[1,2,5]}}
\begin{center}
\begin{tabularx}{\linewidth}{@{}c c X@{}}
\toprule
Amount & Best count & Reason \\
\midrule
0 & 0 & base case: no coins are needed \\
1 & 1 & \(1+dp[0]\) using coin 1 \\
2 & 1 & coin 2 is better than two coin-1 choices \\
5 & 1 & one coin of value 5 \\
6 & 2 & coin 1 plus the best answer for 5 \\
10 & 2 & \(5+5\) \\
11 & 3 & \(1+dp[10]\), producing \(1+5+5\) \\
\bottomrule
\end{tabularx}
\end{center}
\end{dryrunbox}

The sentinel \texttt{amount + 1} is safe because no valid answer can require
more than \texttt{amount} positive-valued coins when coin 1 exists. More
generally, it simply represents a value worse than every possible valid
answer under the problem rules.

\subsection{Code Walkthrough}
The outer loop grows the solved prefix from 1 through \texttt{amount}. For
each coin not larger than the current amount, the code considers taking that
coin last. The remaining state \texttt{current - coin} is smaller and has
already been solved. After all coins are tested, \texttt{dp[current]} is the
minimum among every possible final coin.

\subsection{Edge Cases, Alternatives, and Interview Notes}
Amount zero returns zero. An empty coin list or unreachable amount returns
\(-1\). Duplicate denominations do not change the minimum but perform
unneeded work. A top-down memoized solution uses the same update rule and
is valuable when only a small subset of amounts is reached.

For example, \texttt{coins=[2]} and \texttt{amount=3} leaves \(dp[3]\) at the
sentinel because neither amount 1 nor amount 3 can be formed. This is different
from a valid answer of zero, which occurs only when the requested amount is
zero.

\begin{revisionbox}{How to explain it simply}
I define \(dp[x]\) as the fewest coins needed for amount \(x\). For every
amount, I try each denomination as the final coin and extend the already
best state \(dp[x-c]\). The table prevents repeated recursion and gives
\(O(A m)\) time with \(O(A)\) space.
\end{revisionbox}
\chaptercheckpoint
  {Reusable choices build a target and the objective asks for a minimum.}
  {A DP entry for the minimum coins required by every smaller amount.}
  {Using zero for unreachable states or forgetting the amount-zero base case.}
  {$O(A m)$ time and $O(A)$ space.}
\judgeclassnote


\problemstart{Combination Sum IV}{377}{Medium}{Ordered-count dynamic programming}{One-dimensional DP}

\subsection{The Problem}
Given distinct positive integers and a target, count the ordered sequences that
sum to the target. With \texttt{nums=[1,2,3]} and target \(4\), the answer is
\(7\). For example, \([1,3]\) and \([3,1]\) are different combinations.

\subsection{Reasoning and Loop Order}
Let \(dp[x]\) be the number of ordered sequences totaling \(x\). The empty
sequence is one way to create zero, so \(dp[0]=1\). For every target value
from \(1\) upward, try every number as the final element: \(dp[x] = \sum_{v\le x} dp[x-v].\)
The outer loop must be the current total and the inner loop the choice
number. Reversing them counts unordered coin combinations instead.

\begin{invariantbox}{the table counts every ordered sequence by its final value}
Before computing \(dp[x]\), all counts for smaller totals are complete. The
sequences ending with different final values do not overlap, so adding
\(dp[x-v]\) over every valid \(v\) counts each ordered sequence exactly once.
\end{invariantbox}

\begin{visualbox}{Three last-number choices build the answer for 4}
\centering
\begin{tikzpicture}[font=\small,>={Stealth[length=2.2mm,width=1.55mm]}]
  \node[dsnode,visited] (one) at (0,1.35) {last number 1\\4 ways to make 3};
  \node[dsnode] (two) at (0,0) {last number 2\\2 ways to make 2};
  \node[dsnode,fill=softorange] (three) at (0,-1.35) {last number 3\\1 way to make 1};
  \node[dsnode,active,minimum width=27mm] (answer) at (4.8,0)
    {total ways for 4\\\(4+2+1=7\)};
  \draw[flowarrow,teal] (one.east)--(answer.north west);
  \draw[flowarrow,navy] (two.east)--(answer.west);
  \draw[flowarrow,orange] (three.east)--(answer.south west);
  \foreach \x/\v in {0/1,1/1,2/2,3/4,4/7} {
    \node[arraycell,minimum width=11mm] (dp\x) at ({1.1+1.3*\x},-2.75) {\v};
    \node[below=1mm of dp\x,font=\scriptsize] {\(dp[\x]\)};
  }
  \node[draw=orange,very thick,fit=(dp4),inner sep=1mm] {};
\end{tikzpicture}

\smallskip
For total 4, group the sequences by their last number. The three groups contain
4, 2, and 1 sequence, giving \(4+2+1=7\).
\end{visualbox}

\begin{dryrunbox}{DP for \texttt{[1,2,3]} and target 4}
\begin{center}
\begin{tabular}{c|ccccc}
Target \(x\) & 0 & 1 & 2 & 3 & 4 \\ \hline
\(dp[x]\) & 1 & 1 & 2 & 4 & 7
\end{tabular}
\end{center}
\end{dryrunbox}

\subsection{Pseudocode and C++ Implementation}
\begin{lstlisting}[style=pseudocode]
dp[0] = 1
for total from 1 through target:
    for value in nums:
        if value <= total:
            dp[total] += dp[total - value]
return dp[target]
\end{lstlisting}

\begin{lstlisting}[style=compactcode,caption={Complete C++ solution: Combination Sum IV}]
#include <vector>

using namespace std;

class Solution {
public:
    int combinationSum4(vector<int>& nums, int target) {
        vector<unsigned long long> dp(target + 1, 0);
        dp[0] = 1;
        for (int total = 1; total <= target; ++total) {
            for (const int value : nums) {
                if (value <= total) dp[total] += dp[total - value];
            }
        }
        return static_cast<int>(dp[target]);
    }
};
\end{lstlisting}

\subsection{Why It Works, Complexity, and Common Mistakes}
Every valid sequence totaling \(x\) has one final value \(v\); removing it
leaves a sequence counted by \(dp[x-v]\). These groups do not overlap, so their
counts add exactly. Time is \(O(Tm)\), space is \(O(T)\), where \(T\) is the
target and \(m\) is the number of values. Do not set \(dp[0]=0\), and remember
that order matters in this problem.

\begin{revisionbox}{How to explain it in an interview}
I count ways for every smaller total. The base \(dp[0]=1\) represents choosing
nothing. For each total, I try each number as the last choice and add the ways
to form the remaining total. This produces ordered sequences in \(O(Tm)\).
\end{revisionbox}

\subsection{Why the Loop Order Counts Ordered Sequences}
The worked example contains seven sequences for target 4. The important detail is
that \([1,3]\) and \([3,1]\) are different. The outer loop chooses
the total being built, while the inner loop tries every possible final value.
For each final value, all prefixes that form the remaining total are appended
with that value.

If the loops were reversed---numbers outside and totals inside---the table
would count unordered coin combinations instead. Loop order is part of the
meaning of the state, not merely a coding preference.

\subsection{Complete Dry Run}
\begin{dryrunbox}{Seven sequences for target 4}
\begin{center}
\begin{tabular}{@{}c c l@{}}
\toprule
Total & Ways & Sequences added by their last value \\
\midrule
0 & 1 & empty prefix \\
1 & 1 & \([1]\) \\
2 & 2 & \([1,1],[2]\) \\
3 & 4 & \([1,1,1],[2,1],[1,2],[3]\) \\
4 & 7 & \([1,1,1,1],[2,1,1],[1,2,1],[3,1],\) \\
  &   & \([1,1,2],[2,2],[1,3]\) \\
\bottomrule
\end{tabular}
\end{center}
\end{dryrunbox}

\subsection{Code Walkthrough and Safety}
The unsigned table provides extra intermediate range, while the problem
guarantees the final answer fits a signed 32-bit integer. The base value 1 at
index zero represents the one valid empty prefix. Every access uses
\texttt{total - value} only after confirming that the result is non-negative.

\subsection{Edge Cases and Follow-ups}
Target zero has one empty sequence. A number greater than the target is simply
ignored. Inputs must be positive; allowing zero or negative values would
break the increasing-state dependency and could create infinitely many
sequences. A common follow-up asks for unordered combinations, which requires
reversing the loop order.

\begin{revisionbox}{Practice checkpoint}
Explain why \(dp[0]\) is one even though the answer contains no number. Then
rewrite the update rule for combinations where order does not matter.
\end{revisionbox}
\chaptercheckpoint
  {The problem counts ordered ways to reach a target.}
  {A DP count for each total, beginning with the empty sequence at zero.}
  {Reversing the loops and accidentally counting unordered combinations.}
  {$O(Tm)$ time and $O(T)$ space.}
\judgeclassnote


\patternintoc{Tree Construction}
\problemstart{Construct Binary Tree from Preorder and Inorder Traversal}{105}{Medium}{Recursive partitioning}{Tree and hash map}

\subsection{The Problem}
Preorder lists each root before its subtrees. Inorder lists the left subtree,
then the root, then the right subtree. Reconstruct and return the unique tree
when all values are distinct.

\begin{conceptbox}{Example}
\texttt{preorder=[3,9,20,15,7]} and
\texttt{inorder=[9,3,15,20,7]} reconstruct the tree rooted at \(3\), with
left child \(9\) and right subtree rooted at \(20\).
\end{conceptbox}

\subsection{Recursive Split}
The next unused preorder value is the current root. Its position in inorder
splits the current interval into left and right subtrees. A hash map gives that
position in average \(O(1)\) time, avoiding a repeated linear search.

\begin{invariantbox}{the next preorder value roots the current inorder interval}
Each recursive call owns exactly one consecutive inorder interval. The shared
preorder index points to that interval's root; after the left call consumes
all left-subtree nodes, the right call begins at exactly the next unused
preorder value.
\end{invariantbox}

\begin{visualbox}{Preorder gives the root; inorder shows its left and right sides}
\centering
\begin{tikzpicture}[font=\small]
  \node[dsnode,fill=softorange] (pre) {preorder\\\texttt{3 | 9 | 20 15 7}};
  \node[dsnode,below=10mm of pre] (in) {inorder\\\texttt{9 | 3 | 15 20 7}};
  \draw[flowarrow,orange] (pre.south)--node[edgelabel,right,font=\scriptsize]{root 3} (in.north);
  \node[below left=9mm and 13mm of in,text=teal] {left side: \texttt{[9]}};
  \node[below right=9mm and 13mm of in,text=navy] {right side: \texttt{[15,20,7]}};
\end{tikzpicture}
\end{visualbox}

\begin{visualbox}{The map finds each root; the stack remembers unfinished ranges}
\centering
\begin{tikzpicture}[font=\small,>={Stealth[length=2.2mm,width=1.55mm]}]
  \node[dsnode,align=left,fill=softteal] (map) {
    \textbf{position of each value}\\
    \(9\to0\)\quad \(3\to1\)\\
    \(15\to2\)\quad \(20\to3\)\\
    \(7\to4\)};
  \node[draw=navy,fill=softblue,minimum width=46mm,minimum height=8mm,
    right=28mm of map] (b0) {whole range: root 3};
  \node[draw=navy,fill=softblue,minimum width=46mm,minimum height=8mm,
    above=0mm of b0] (b1) {left range: root 9};
  \node[draw=orange,fill=softorange,minimum width=46mm,minimum height=8mm,
    above=0mm of b1] (b2) {empty range: stop};
  \node[above=5mm of map,text=teal,font=\bfseries] {find the root quickly};
  \node[right=7mm of b2,text=orange,font=\bfseries,align=left] {working\\now};
  \draw[flowarrow,teal] (map.east)--node[edgelabel,above,font=\scriptsize]{root position} (b0.west);
\end{tikzpicture}

\smallskip
The map tells us where the root sits. The stack remembers which left and right
parts still need to be built.
\end{visualbox}

\subsection{Pseudocode}
\begin{lstlisting}[style=pseudocode]
build(left, right):
    if left > right: return null
    rootValue = preorder[preorderIndex++]
    root = new TreeNode(rootValue)
    middle = inorderPosition[rootValue]
    root.left = build(left, middle - 1)
    root.right = build(middle + 1, right)
    return root
\end{lstlisting}

\subsection{C++ Implementation}
\begin{lstlisting}[style=compactcode,caption={Complete C++ solution: Construct Binary Tree}]
#include <unordered_map>
#include <vector>

using namespace std;

class Solution {
    const vector<int>* preorderValues = nullptr;
    unordered_map<int, int> inorderIndex;
    int preorderIndex = 0;

    TreeNode* build(int left, int right) {
        if (left > right) return nullptr;
        const int rootValue = (*preorderValues)[preorderIndex++];
        TreeNode* root = new TreeNode(rootValue);
        const int middle = inorderIndex[rootValue];
        root->left = build(left, middle - 1);
        root->right = build(middle + 1, right);
        return root;
    }

public:
    TreeNode* buildTree(vector<int>& preorder, vector<int>& inorder) {
        preorderValues = &preorder;
        preorderIndex = 0;
        inorderIndex.clear();
        for (int i = 0; i < static_cast<int>(inorder.size()); ++i) {
            inorderIndex[inorder[i]] = i;
        }
        return build(0, static_cast<int>(inorder.size()) - 1);
    }
};
\end{lstlisting}
\judgeclassnote

\subsection{Why It Works, Time, and Space}
Each call consumes exactly one preorder root and assigns it the inorder range
that contains exactly its descendants. Building the left range before the
right range matches preorder order. Every node is created once, so time and
stored map space are \(O(n)\); recursion uses \(O(h)\) stack space.

\begin{revisionbox}{Common mistakes}
Build the left subtree before the right subtree because preorder lists the
entire left subtree first. Use interval boundaries instead of copying array
slices. The distinct-value assumption is what makes the inorder index unique.
\end{revisionbox}

\subsection{Direct Method and Its Limitation}
A simple recursive version can search the full inorder array for every new
root. It constructs the correct tree, but the repeated linear searches cost
\(O(n^2)\) for a skewed tree. Building a value-to-index map once changes every
root lookup to average \(O(1)\), producing an \(O(n)\) construction.

\subsection{Detailed Recursive Walkthrough}
\begin{dryrunbox}{Preorder \texttt{[3,9,20,15,7]}, inorder \texttt{[9,3,15,20,7]}}
\begin{center}
\begin{tabularx}{\linewidth}{@{}c c c X@{}}
\toprule
Call & Next preorder root & Inorder interval & Result \\
\midrule
1 & 3 & \([0,4]\) & split at index 1; left \([0,0]\), right \([2,4]\) \\
2 & 9 & \([0,0]\) & leaf; both child intervals are empty \\
3 & 20 & \([2,4]\) & split at index 3 \\
4 & 15 & \([2,2]\) & left leaf of 20 \\
5 & 7 & \([4,4]\) & right leaf of 20 \\
\bottomrule
\end{tabularx}
\end{center}
\end{dryrunbox}

\begin{visualbox}{The tree after all left and right parts are built}
\centering
\begin{tikzpicture}[font=\small,>={Stealth[length=2.2mm,width=1.55mm]}]
  \node[trienode,active] (r) {3};
  \node[trienode,below left=13mm and 22mm of r] (n9) {9};
  \node[trienode,below right=13mm and 22mm of r] (n20) {20};
  \node[trienode,visited,below left=13mm and 9mm of n20] (n15) {15};
  \node[trienode,visited,below right=13mm and 9mm of n20] (n7) {7};
  \draw[flowarrow] (r)--node[edgelabel,above left,font=\scriptsize]{left side} (n9);
  \draw[flowarrow] (r)--node[edgelabel,above right,font=\scriptsize]{right side} (n20);
  \draw[flowarrow] (n20)--(n15); \draw[flowarrow] (n20)--(n7);
  \node[below=5mm of n9,font=\scriptsize,text=teal]{position 0};
  \node[below=5mm of n15,font=\scriptsize,text=teal]{position 2};
  \node[below=5mm of n7,font=\scriptsize,text=teal]{position 4};
\end{tikzpicture}
\end{visualbox}

The preorder pointer advances exactly once per created node. The inorder
boundaries determine when a subtree is empty; no array slices need to be
allocated.

\subsection{Code Walkthrough}
The public method stores a pointer to the preorder vector, resets the shared
index, and creates the inorder lookup map. The helper reads the next root,
allocates it, finds its inorder position, and recursively builds the left
interval before the right interval. The base case \texttt{left > right}
returns a null child.

\subsection{Assumptions, Edge Cases, and Follow-ups}
The traversals must contain the same distinct values and represent a valid
tree. Empty arrays produce a null root. A skewed tree uses \(O(n)\) recursion
depth. Duplicate values make the reconstruction ambiguous unless extra
occurrence information is supplied. A useful variation reconstructs from
inorder and postorder; in that case roots are consumed from the end and the
right subtree must be built first.

\begin{revisionbox}{How to explain it simply}
Preorder tells me the next root. Its position in inorder divides the exact
left and right descendant ranges. A hash map removes repeated searches, and a
single preorder index ensures each node is consumed once, for \(O(n)\) time.
\end{revisionbox}
\chaptercheckpoint
  {Two traversals reveal roots and subtree boundaries in complementary orders.}
  {A preorder cursor, inorder index map, and recursive interval boundaries.}
  {Building the right subtree first or copying slices at every call.}
  {$O(n)$ time, $O(n)$ map space, and $O(h)$ recursion space.}